\documentclass[aspectratio=169,11pt]{beamer}
\usetheme{metropolis}
\usepackage{booktabs}
\usepackage{amsmath,amssymb}
\usepackage{graphicx}
\usepackage{xcolor}

\definecolor{accent}{HTML}{2E86AB}
\definecolor{goodgreen}{HTML}{2E7D32}
\definecolor{warnred}{HTML}{C62828}
\definecolor{muted}{HTML}{757575}

\setbeamercolor{frametitle}{bg=accent,fg=white}
\setbeamercolor{progress bar}{fg=accent}

\title{Language Features under Weight Quantization}
\subtitle{Activation Agreement and Intervention Effects}
\author{Hemang Joshi \and Eashaan Thakur \and Shardul Joshi \and Nikhita Ravi}
\institute{IIIT Hyderabad -- ANLP Monsoon 2026}
\date{}

\begin{document}

\begin{frame}[plain]
\titlepage
\end{frame}

% ──────────────────────────────────────────────
\begin{frame}{The question}
\Large
When a model's weights are rounded to fewer bits,\\[6pt]
do the internal features that \textbf{explain} its behaviour\\[6pt]
still work the same way?

\vspace{12pt}
\normalsize
\textcolor{muted}{We study language-selective features across 27 languages,\\
three quantizers, and two Gemma model sizes.}
\end{frame}

% ──────────────────────────────────────────────
\begin{frame}{Why this matters}
\begin{itemize}
\item Quantization makes LLMs cheaper to run, but can change behaviour
\item Effects differ across languages -- low-resource languages may suffer more
\item Sparse autoencoders (SAEs) find features that fire for specific languages
\item Prior work (Duan, 2026) measured \emph{activation survival} under quantization
\end{itemize}

\vspace{8pt}
\textbf{Our contribution:} we go beyond activation survival and ask whether features still \emph{do the same thing} when edited
\end{frame}

% ──────────────────────────────────────────────
\begin{frame}{Three things that come apart}
\centering
\begin{tabular}{lp{7cm}c}
\toprule
\textbf{Measure} & \textbf{What it tests} & \textbf{INT4 legacy}\\
\midrule
Activation survival ($R_1 > 0.9$) & Does the feature fire the same way? & \textcolor{warnred}{13.2\%}\\[4pt]
Functional effect ($E$) & Does removing it still hurt the language? & \textcolor{goodgreen}{27/27}\\[4pt]
Effect identity ($W_1$) & Does it still do the same thing to output? & \textcolor{goodgreen}{99.7\%}\\
\bottomrule
\end{tabular}

\vspace{12pt}
The activation cutoff says 87\% of features are ``dead.''\\
But 99.7\% can still be identified from what they do.
\end{frame}

% ──────────────────────────────────────────────
\begin{frame}{Setup}
\begin{columns}[T]
\column{0.5\textwidth}
\textbf{Models}
\begin{itemize}
\item Gemma-2-2B (main)
\item Gemma-2-9B (confirmation)
\item Llama-3.2-1B (cross-check)
\end{itemize}

\vspace{6pt}
\textbf{Languages}\\
27 from FLORES-200\\
9 high / 9 mid / 9 low resource

\column{0.5\textwidth}
\textbf{Quantizers}
\begin{itemize}
\item Legacy RTN (from prior work)
\item Matched RTN (same method, better implementation)
\item GPTQ (calibrated rounding)
\end{itemize}

\vspace{6pt}
\textbf{SAE}\\
Gemma Scope JumpReLU, 16k features\\
One fixed layer per model size
\end{columns}
\end{frame}

% ──────────────────────────────────────────────
\begin{frame}{What we measure}
\textbf{For each language}, select 64 features that fire most selectively for it

\vspace{6pt}
\begin{enumerate}
\item \textbf{Zero-ablate} the group: set those 64 features to zero
\item Measure the increase in next-token loss (NLL)
\item Compare against \textbf{matched controls}: non-selective features with similar firing rates
\item Net harm $E$ = group harm $-$ control harm
\end{enumerate}

\vspace{8pt}
Do the same at full precision and at each bit width.\\
The \textbf{retention ratio} $\rho = E^{(b)} / E^{\text{FP}}$ tells us if the edit matters more or less after quantization.
\end{frame}

% ──────────────────────────────────────────────
\begin{frame}{Correction: the missing BOS token}
Reproducing Duan (2026), we found a methodological issue:

\vspace{6pt}
\begin{center}
\begin{tabular}{lcc}
\toprule
& \textbf{Without BOS} & \textbf{With BOS}\\
\midrule
FP perplexity & 458.4 & \textbf{16.1}\\
INT6 survival & 51\% & \textbf{97.5\%}\\
INT7 survival & 81\% & \textbf{$>$99.9\%}\\
\bottomrule
\end{tabular}
\end{center}

\vspace{6pt}
One missing token at the start of the input made the model look 28$\times$ worse at prediction and made half the features look ``dead'' at INT6.

\vspace{4pt}
\textcolor{muted}{We corrected our protocol and retired our own first pipeline after finding separate flaws.}
\end{frame}

% ──────────────────────────────────────────────
\begin{frame}{Features matter in all 27 languages}

Removing a language's feature group hurts next-token prediction more than removing matched controls.

\vspace{6pt}
\begin{itemize}
\item \textbf{FP:} positive excess harm in 27/27 languages (median 0.262 nats/token)
\item \textbf{INT6:} retention ratio $\approx$ 1.0 (features matter about as much)
\item \textbf{INT4:} retention ratio 1.54 (legacy RTN) -- the edit hurts \emph{more}
\end{itemize}

\vspace{8pt}
Gradual suppression confirms dose--response in all 27 languages: as you turn features down, harm increases monotonically.
\end{frame}

% ──────────────────────────────────────────────
\begin{frame}{But features carry almost none of the damage}

If features are hurt more at INT4, maybe fixing them fixes the model?

\vspace{8pt}
\textbf{Rescue experiment:} patch FP activations back into the quantized model.

\vspace{6pt}
\begin{center}
\begin{tabular}{lc}
\toprule
What we restore & \% of INT4 loss gap recovered\\
\midrule
Language group (64 features) & \textcolor{warnred}{2.6\%}\\
Full residual stream & 39.3\%\\
\bottomrule
\end{tabular}
\end{center}

\vspace{6pt}
The amplified ratio does not mean the features carry the quantization damage.
The compressed model is just more sensitive to any perturbation.
\end{frame}

% ──────────────────────────────────────────────
\begin{frame}{The direction-swap control}

Is the INT4 similarity loss \emph{specific} to each feature?

\vspace{6pt}
\textbf{Control:} ablate a \emph{different} feature's direction at the same positions and strengths.

\vspace{6pt}
\begin{center}
\begin{tabular}{lcc}
\toprule
& Real feature & Swapped direction\\
\midrule
Median $W_1$ at INT4 & 0.842 & 0.828\\
\bottomrule
\end{tabular}
\end{center}

\vspace{6pt}
The gap is 0.014 -- consistent at every bit width.\\[4pt]
The similarity loss reflects how the \emph{downstream computation} drifts under quantization for \textbf{any} perturbation, not anything specific to the feature.
\end{frame}

% ──────────────────────────────────────────────
\begin{frame}{GPTQ helps substantially}
\centering
\textbf{Three quantizers on Gemma-2-2B at INT4}

\vspace{8pt}
\begin{tabular}{lccc}
\toprule
& Legacy RTN & Matched RTN & GPTQ\\
\midrule
Survival (\%) & 12.7 & 18.6 & \textbf{30.5}\\
Mean loss (nats/tok) & 0.590 & 0.385 & \textbf{0.130}\\
Median $\rho$ & 1.54 & 1.36 & \textbf{1.22}\\
Median $W_1$ & 0.839 & 0.875 & \textbf{0.923}\\
\bottomrule
\end{tabular}

\vspace{10pt}
\begin{itemize}
\item GPTQ roughly halves the INT4 amplification
\item Matched RTN alone (better implementation, same method) closes about a third of the gap
\item At INT6, all three quantizers preserve features well
\end{itemize}
\end{frame}

% ──────────────────────────────────────────────
\begin{frame}{Gemma-2-9B: confirmed}

We ran the full study on a second model size with all three quantizers.

\vspace{8pt}
\centering
\begin{tabular}{lccc}
\toprule
\textbf{9B at INT4} & Legacy & Matched & GPTQ\\
\midrule
Survival (\%) & 22.3 & 32.9 & \textbf{46.6}\\
Loss (nats/tok) & 0.410 & 0.287 & \textbf{0.088}\\
Amplification $\rho$ & 1.69 & 1.49 & \textbf{1.17}\\
\bottomrule
\end{tabular}

\vspace{10pt}
\begin{itemize}
\item GPTQ below legacy in \textbf{26 of 27} languages
\item All frozen predictions (written before any 9B data existed) were met
\item The pattern holds across model sizes within the Gemma family
\end{itemize}
\end{frame}

% ──────────────────────────────────────────────
\begin{frame}{Honest limits}
\begin{columns}[T]
\column{0.5\textwidth}
\textbf{Belebele null}\\[4pt]
Our only confirmatory test.\\
Removing a language's features does \emph{not} change reading-comprehension accuracy.\\[4pt]
FP effect: $-0.0001$ nats/item\\
$[-0.0035, +0.0031]$\\[4pt]
The features matter for text prediction but not for this comprehension task.

\column{0.5\textwidth}
\textbf{Llama-3.2-1B: not consistent}\\[4pt]
Positive FP effect in only 21/27 languages (rule: 24).\\
INT4 ratio: 0.86 (no amplification).\\[4pt]
But model family, SAE architecture, and tokenizer all change together, so we cannot isolate the cause.
\end{columns}

\vspace{8pt}
\textcolor{muted}{Most FLORES analyses are exploratory (sentences reused). Only Belebele was held out.}
\end{frame}

% ──────────────────────────────────────────────
\begin{frame}{What we claim}
\begin{enumerate}
\item Language-selective SAE features contribute to text prediction in \textbf{all 27 languages} on both Gemma sizes
\item \textbf{Activation survival understates preservation}: 99.7\% of features retain their identity even when $R_1$ says they are dead
\item The INT4 similarity loss is \textbf{generic}, not feature-specific (direction-swap control)
\item The features carry only \textbf{2.6\%} of INT4 quantization damage (rescue study)
\item \textbf{GPTQ roughly halves} the INT4 amplification compared to RTN
\item The pattern \textbf{holds on 9B} but \textbf{not on Llama}
\end{enumerate}

\vspace{6pt}
\textcolor{muted}{The Belebele null and rescue limit show where the features \emph{don't} matter.}
\end{frame}

% ──────────────────────────────────────────────
\begin{frame}{Corrections and integrity}
\begin{itemize}
\item Found and reported the BOS error in prior work (reproduced their numbers first)
\item Retired our own first pipeline after finding separate flaws
\item Two 9B predictions were defective (reversed inequality, tautological criterion) -- documented, not counted
\item Frozen predictions written before data; code hashes locked; controls pre-specified
\item Belebele null and Llama failure reported as results, not hidden
\end{itemize}

\vspace{6pt}
\textbf{Resources:}\\
Code: \texttt{github.com/TH4KUR/always-late}\\
Artifacts: \texttt{huggingface.co/datasets/karma-skz/always-late-artifacts}\\
Runs: \texttt{wandb.ai/karma-skz-/anlp-project}
\end{frame}

\end{document}
